\chapter{Estudio de la tecnología}

\begin{quotation} [Físico] {Stephen Hawking (1942--2018)}
El mayor enemigo del conocimiento no es la ignorancia, es la ilusión del conocimiento.
\end{quotation}

\begin{abstract}
Vamos a explicar las diferentes tecnologías que hemos utilizado para el desarrollo de la aplicación.
\end{abstract}

\section{Introducción}

En este capítulo vamos a realizar un estudio de las distintas tecnologías que se han utilizado para desarrollar la aplicación web.

En el primer punto se tratará de la base de datos empleada, la cual ha sido PostgreSQL, una base de datos relacional orientada a objetos y de código abierto.

En el segundo punto se hablará de PM2 y NGINX. PM2 es un gestor de procesos de producción para aplicaciones NodeJS, el cual se utilizará para mantener las aplicaciones en ejecución, gestionar el reinicio automático y monitorizar el rendimiento. NGINX es un servidor web y proxy inverso que se utilizará para gestionar el tráfico de red y mejorar el rendimiento y la seguridad de la aplicación.

En el tercer punto se tratará de NodeJS y ExpressJS, los cuales se han empleado para la implementación del servidor (backend). NodeJS es un entorno de ejecución de JavaScript, mientras que ExpressJS es un framework minimalista y flexible para aplicaciones web en NodeJS.

En el cuarto punto se tratará sobre React, una biblioteca Javascript de código abierto con la finalidad de facilitar el desarrollo de aplicaciones en una sola página que se ha empleado para la implementación del cliente.

En el último punto se tratará sobre Vite y CSS. Vite es una herramienta de frontend que ayudará a crear el proyecto y que tanto el desarrollo como la construcción final sea lo más sencilla posible. El uso de CSS se ha elegido para no sobrecargar el proyecto con librerías o frameworks de estilos adicionales.

\newpage

\section{PostgreSQL}

\begin{figure}[h]
\centering
 \caption{Tecnología - PostgreSQL}
\end{figure}

PostgreSQL es la base de datos de código abierto más avanzada del mundo y es un poderoso sistema de base de datos objeto-relacional de código abierto que utiliza y extiende el lenguaje SQL combinado con muchas características que almacenan y escalan de manera segura las cargas de trabajo de datos más complicadas\footnote{Definición de PostgreSQL (\href{https://hoplasoftware.com/postgresql-base-de-datos-relacional-codigo-abierto/}{en la siguiente referencia})}.\\

PostgreSQL se ha ganado una sólida reputación por su arquitectura comprobada, confiabilidad, integridad de datos, conjunto sólido de funciones, extensibilidad y la dedicación de la comunidad de código abierto detrás del software para ofrecer soluciones innovadoras y de alto rendimiento. PostgreSQL se ejecuta en todos los principales sistemas operativos, cumple con ACID desde 2001 \footnote{Características de PostgreSQL (\href{https://www.postgresql.org/about/}{en la siguiente referencia})}.\\

\begin{figure}[h]
\centering
    	\caption{Tecnología - Principios ACID}
\end{figure}

\newpage

\section{PM2 y NGINX}
En este punto vamos a profundizar en el uso de PM2 y NGINX, dos herramientas clave para la gestión y optimización del servidor donde se albergará la aplicación.

\begin{figure}[h]
  \centering
   \caption{Tecnología - PM2}
\end{figure}
  
  PM2 es un gestor de procesos de producción para aplicaciones NodeJS. Esta herramienta ofrece una serie de características que facilitan la administración y el rendimiento de las aplicaciones en un entorno de producción:

\begin{itemize}

  \item \textbf{Manejo de Procesos}: PM2 permite ejecutar múltiples instancias de una aplicación NodeJS, maximizando el uso de los recursos del servidor y mejorando la capacidad de manejo de cargas de trabajo.
  
  \item \textbf{Reinicio Automático}: En caso de fallos o caídas de la aplicación, PM2 reinicia automáticamente la instancia para asegurar la continuidad del servicio sin necesidad de intervención manual.
  
  \item \textbf{Monitorización y Logs}: PM2 proporciona herramientas de monitorización en tiempo real y gestión de logs, lo que facilita la identificación y resolución de problemas. Esto incluye estadísticas de CPU, memoria y tiempo de respuesta.
  
  \item \textbf{Cluster Mode}: PM2 permite ejecutar la aplicación en modo clúster, distribuyendo las solicitudes entrantes entre las diversas instancias de NodeJS para mejorar el rendimiento y la escalabilidad.
  
  \item \textbf{Despliegue}:PM2 simplifica el proceso de despliegue de aplicaciones, permitiendo el despliegue continuo y la actualización sin tiempo de inactividad.

\end{itemize}

\newpage
\begin{figure}[h]
   \caption{Tecnología - NGINX}
\end{figure}
  
NGINX es un servidor web y proxy inverso que se utilizará para mejorar el rendimiento y la seguridad de la aplicación. Algunas de las características más relevantes de NGINX son:

\begin{itemize}

  \item \textbf{Proxy Inverso}: NGINX actúa como intermediario entre los usuarios y el servidor de la aplicación, manejando las solicitudes entrantes y distribuyéndolas de manera eficiente. Esto ayuda a balancear la carga de trabajo y mejorar la velocidad de respuesta.
  
  \item \textbf{Seguridad}: NGINX proporciona características avanzadas de seguridad, como la gestión de certificados SSL/TLS, para asegurar las comunicaciones entre los usuarios y el servidor. Además, puede proteger la aplicación de ataques comunes como DDoS, limitando el número de solicitudes por usuario.
  
  \item \textbf{Caché de Contenido}: NGINX puede almacenar en caché el contenido estático y dinámico, reduciendo la carga en el servidor y acelerando la entrega de contenido a los usuarios
  
  \item \textbf{Balanceo de Carga}: NGINX distribuye el tráfico de manera equilibrada entre las diferentes instancias de la aplicación, mejorando la disponibilidad y la capacidad de respuesta del servicio.
  
  \item \textbf{Compresión y Optimización}: NGINX puede comprimir el contenido antes de enviarlo a los usuarios, reduciendo el uso de ancho de banda y acelerando el tiempo de carga de las páginas.

\end{itemize}

El uso combinado de PM2 y NGINX proporciona una solución robusta y escalable para gestionar y optimizar el servidor de la aplicación, asegurando tanto el rendimiento como la seguridad del servicio.

\newpage

\section{NodeJS y ExpressJS}

En este punto, profundizaremos en el uso de NodeJS y ExpressJS, dos tecnologías clave para la implementación del servidor (backend) de nuestra aplicación web.

\begin{figure}[h]
\centering
 \caption{Tecnología - NodeJS}
\end{figure}

NodeJS es un entorno de ejecución de JavaScript basado en el motor V8 de Google Chrome. Permite ejecutar código JavaScript en el servidor, lo que hace posible desarrollar tanto el frontend como el backend de una aplicación utilizando el mismo lenguaje de programación. Algunas de las características más destacadas de NodeJS son:

\begin{itemize}

\item \textbf{Asincronía y Event-Driven}: NodeJS utiliza un modelo de programación asincrónico y basado en eventos, lo que le permite manejar múltiples conexiones concurrentes sin bloquear el hilo principal. Esto se traduce en un alto rendimiento y una escalabilidad eficiente.

\item \textbf{NPM (Node Package Manager)}: NodeJS viene acompañado de NPM, el gestor de paquetes más grande del mundo. Con NPM, los desarrolladores pueden acceder a miles de bibliotecas y módulos reutilizables que facilitan el desarrollo de aplicaciones.

\item \textbf{Cross-Platform}: NodeJS es compatible con múltiples sistemas operativos, incluidos Linux, Windows y macOS, lo que permite desarrollar y desplegar aplicaciones en diferentes entornos sin problemas.

\item \textbf{Comunidad Activa}: La comunidad de NodeJS es grande y activa, lo que significa que hay abundantes recursos, tutoriales y soporte disponibles para los desarrolladores.

\end{itemize}

\newpage

\begin{figure}[h]
\centering
 \caption{Tecnología - ExpressJS}
\end{figure}

ExpressJS es un framework minimalista y flexible para aplicaciones web en NodeJS. Se ha convertido en una de las herramientas más populares para el desarrollo de servidores debido a su simplicidad y capacidad de extensión. Algunas de las características más importantes de ExpressJS incluyen:

\begin{itemize}

  \item \textbf{Rutas Simples y Flexibles}: ExpressJS proporciona una sintaxis sencilla para definir rutas y manejar solicitudes HTTP. Esto facilita la creación de API RESTful y la gestión de diferentes endpoints en la aplicación.
  
  \item \textbf{Middleware}: ExpressJS utiliza un sistema de middleware que permite añadir funcionalidades adicionales a las solicitudes y respuestas de HTTP. Los middlewares se pueden encadenar para realizar tareas como autenticación, gestión de sesiones, validación de datos y más.
  
  \item \textbf{Integración con Bases de Datos}: ExpressJS se integra fácilmente con diversas bases de datos, incluyendo MongoDB, MySQL y PostgreSQL, lo que facilita la gestión y acceso a datos desde el servidor.
  
  \item \textbf{Extensible y Modular}: La arquitectura modular de ExpressJS permite a los desarrolladores añadir y eliminar funcionalidades según las necesidades del proyecto, sin afectar el rendimiento o la estructura básica de la aplicación.
  
\end{itemize}

El uso combinado de NodeJS y ExpressJS proporciona una solución robusta y eficiente para la implementación del backend de la aplicación web. NodeJS aporta el entorno de ejecución y las capacidades asincrónicas, mientras que ExpressJS ofrece las herramientas necesarias para gestionar rutas, middleware y la interacción con bases de datos.

\newpage

\section{React}

\begin{figure}[h]
\centering
 \caption{Tecnología - React}
\end{figure}

React es una biblioteca Javascript de código abierto diseñada para crear interfaces de usuario con el objetivo de facilitar el desarrollo de aplicaciones en una sola página. Es mantenido por Facebook y la comunidad de software libre \footnote{Definición de React. (\href{https://es.wikipedia.org/wiki/React}{en la siguiente referencia})}. \\

React se presenta como una alternativa ideal para hacer todo tipo de aplicaciones web, \textbf{SPA (Single Page Application}. Esto se debe al contar con un completo ecosistema de módulos, herramientas y componentes que posibilitan desarrollar funcionalidades complejas en poco tiempo.\\

Con esta tecnología solo desarrollamos el \textbf{frontend}, por lo que esta tecnología esta integrada junto con Spring mencionado anteriormente para desarrollar ambas partes de la aplicación, por eso podemos decir que React es la "V" si seguimos el patrón de diseño de \textbf{MVC (Modelo-Vista-Controlador)}. \\

El elemento característico de esta librería se identifica como el componente, una pieza central de la interfaz de usuario. Debemos tener en cuenta que, al diseñar con React, se crean componentes independientes y reusables para crear interfaces modulares y adaptables. Todas las aplicaciones desarrolladas con React, tienen un componente al que siempre se le denomina “raíz” y, a su vez, tiene otros llamados “hijos” quienes conforman un árbol de componentes que están entrelazados entre sí \footnote{Elemento característico de React. (\href{https://www.coderhouse.es/blog/que-es-react-js}{en la siguiente referencia})}.\cite{react}

\newpage

Los componentes que he desarrollado en mi aplicación tienen extensión \textbf{JSX}, el cual se trata de una extensión de la sintaxis de JavaScript que nos permite compilar un objeto JavaScript para mapear un elemento del DOM (Objeto de Documento). A través del JSX, React crea una representación en memoria del DOM denominada Virtual DOM que, a diferencia de la común, pesa muy poco y usa menos recursos para ser creado. \\

Con esto, si un componente cambia su estado y la información a ser renderizada, React lo que hace es comparar los cambios realizados en el Virtual DOM, observa cuáles elementos se han modificado y actualiza solo esa parte en el DOM real. En consecuencia a este funcionamiento, la aplicación se ejecuta más rápidamente. El DOM Virtual supervisa los cambios, cuando ocurren, se comparan con la interfaz original y se aplican las modificaciones solo donde sea necesario. \\

\textbf{¿Qué ventajas nos ofrece React?}

\begin{itemize}

\item Ofrece facilidades para el desarrollo de app orientadas al frontend.
\item Destaca por su rendimiento, flexibilidad y la organización del código.
\item Es capaz de generar el DOM de forma dinámica, evitando que se renderice toda la página completa.
\item Garantiza una mejor experiencia de usuario, mayor fluidez en las interfaces e interactividad.
\item Posee una arquitectura de desarrollo extendible, lo que permite la encapsulación del código en componentes.

\end{itemize}

\textbf{¿Por qué React?}\\
La motivación principal de elegir React para la realización del cliente en la aplicación, fue el poder ampliar mis conocimientos en una tecnología dedicada al frontend y con un lenguaje de programación poco usado en el master, ya que nunca había estudiado el lenguaje de programación JavaScript.\\
Con todo lo mencionado anteriormente, vi como una oportunidad en este proyecto el poder realizar una aplicación empleando esta tecnología debido al gran auge y el apoyo de grandes compañías.\\

\newpage

\section{Vite}

\begin{figure}[h]
\centering
 \caption{Tecnología - Vite}
\end{figure}

Vite es una herramienta de compilación que tiene como objetivo proporcionar una experiencia de desarrollo más rápida y ágil para proyectos web modernos\footnote{Definición Vite. (\href{https://es.vitejs.dev/guide/}{en la siguiente referencia})}. Consta de dos partes principales:

\begin{itemize}

\item Un servidor de desarrollo que proporciona mejoras enriquecidas de funcionalidades sobre módulos ES nativos, por ejemplo Hot Module Replacement (HMR) extremadamente rápido.
\item Un comando de compilación que empaqueta tu código con Rollup, preconfigurado para generar recursos estáticos altamente optimizados para producción.

\end{itemize}

Actualmente, Vite soporta la creación tanto de proyectos vanilla (sin utilizar frameworks), como proyectos utilizando Vue, React, Preact, Svelte o Lit (tanto usando Javascript como Typescript).\\

\textbf{Crear un proyecto con Vite} \\

Para empezar, necesitamos tener instalado \textbf{NodeJS}. A continuación, abriremos una terminal y utilizaremos el comando npm para comenzar a crear el proyecto.

En la terminal escribimos el comando npm create vite, en este comando podemos añadir el nombre de la carpeta y se encargará de crear la carpeta del proyecto. \\

A continuación se instalará el paquete de create-vite, el cual es el encargado de crear la aplicación web. Al instalar el paquete, nos aparecerá las distintas tecnologías que permite esta herramienta. En nuestro caso, seleccionamos la plantilla de react, el comando es de la forma npm create tfg -- --template react. 

El siguiente comando a ejecutar dentor de la carpeta del proyecto será npm install. Después de ejecutar el comando, ya tendremos listo para arrancarlo mediante el comando npm run dev y comenzar a desarrollar. 

En el carpeta del proyecto encontraremos un fichero de vite.config.js. Se trata de un fichero de configuración donde se puede establecer algunas configuraciones para el correcto funcionamiento del empaquetador en el proyecto que estamos desarrollando.

\newpage

\section{CSS}

\begin{figure}[h]
  \centering
   \caption{Tecnología - CSS}
  \end{figure}

Para este proyecto hemos optado por utilizar CSS sin ningún framework o librería adicional.\\

\textbf{¿Por qué CSS?}\\

\begin{itemize}

  \item \textbf{Control Completo y Personalización Total}: Permite tener control absoluto sobre cada aspecto del diseño de la aplicación. Esto significa que se puede personalizar hasta el más mínimo detalle según las necesidades. No estás limitado por las restricciones o reglas predefinidas que una librería o framework podría imponer.
  
  \item \textbf{Mejor Rendimiento}: El sitio web no necesita cargar archivos adicionales como los que vienen con las librerías y frameworks. Esto puede resultar en tiempos de carga más rápidos, lo cual es crucial para la experiencia del usuario. Las librerías suelen incluir muchas funcionalidades que quizás no son necesarias para este proyecto, lo que puede aumentar innecesariamente el tamaño del mismo.
  
  \item \textbf{Flexibilidad y Adaptabilidad}: Ofrece la flexibilidad de utilizar cualquier técnica o metodología que se considere adecuada al proyecto. Esto te permite adaptar el diseño a las necesidades específicas del proyecto sin estar atado a una estructura rígida.
  
  \item \textbf{Compatibilidad}: CSS es soportado de manera nativa por todos los navegadores modernos, lo que elimina la necesidad de configuraciones adicionales para asegurar que la aplicación funcione correctamente en diferentes entornos. Además, no dependes de actualizaciones de librerías que podrían potencialmente romper tu código o introducir inconsistencias.
  
  \item \textbf{Menos Dependencias}: Al no depender de librerías o frameworks externos, se reduce el riesgo de que una actualización externa afecte al proyecto.Además, nos permite tener un control total sobre el mantenimiento y las actualizaciones del propio código, lo que puede ser más seguro y estable a largo plazo.

  \end{itemize}